% Guia do professor — Capítulo 21: Processos Climáticos Naturais
\section{Capítulo 21 --- Processos Climáticos Naturais}

\subsection*{Visão geral e ritmo}

A ponte tempo$\to$clima: feedbacks (a álgebra do amplificador),
escalas lentas e variabilidade interna. Quatro modelos de brinquedo
carregam tudo — gelo-albedo, Milankovitch, Daisyworld e o oscilador
ENSO — e cada um é um clássico da literatura rodando no notebook.
\textbf{Ritmo: 2 aulas de 2\,h.}

\begin{itemize}
  \item \textbf{Aula 1} — tempo × clima (por que a estatística é
        previsível), álgebra de feedbacks, bifurcação gelo-albedo
        (Caso~1 — o gráfico da aula), Milankovitch (Caso~2).
  \item \textbf{Aula 2} — Daisyworld (Caso~3), variabilidade interna
        e ENSO (Caso~4), vulcões e Sol, fechamento para o módulo extra (Cap.~23).
\end{itemize}

\subsection*{Objetivos de aprendizagem}

\begin{enumerate}
  \item usar a álgebra de feedbacks ($\Delta T =
        \lambda_0\Delta F/(1-f)$) e classificar os feedbacks
        principais;
  \item explicar múltiplos equilíbrios e histerese do gelo-albedo;
  \item associar os três ciclos orbitais a seus períodos e ao gatilho
        de 65°N;
  \item interpretar o Daisyworld como feedback negativo emergente;
  \item descrever o mecanismo do ENSO e sua previsibilidade sazonal.
\end{enumerate}

\subsection*{Pré-requisitos a reativar}

Balanço radiativo e $T_e$ (Cap.~2); Clausius--Clapeyron para o
feedback do vapor (Cap.~4); insolação orbital (Cap.~2, Caso~2);
alísios e Walker (Cap.~11); caos vs.\ estatística (Cap.~20);
realimentação (eletrônica básica).

\subsection*{Armadilhas conceituais frequentes}

\begin{itemize}
  \item \textbf{``O clima sempre mudou, logo\ldots''}: sim — e cada
        mudança teve CAUSA física identificável (órbita, vulcões,
        CO$_2$); entender as naturais é o que permite atribuir a
        atual (Cap.~23).
  \item \textbf{Feedback positivo $=$ explosão}: com $f<1$ é
        amplificação finita; a fuga só ocorre em $f\to1$ (bola de
        neve).
  \item \textbf{Milankovitch como ``causa direta''}: é o metrônomo;
        a amplitude vem dos feedbacks (gelo, CO$_2$) — N2 mostra a
        retificação.
  \item \textbf{Daisyworld como Gaia teleológica}: não há intenção —
        só seleção $+$ acoplamento; o modelo existe exatamente para
        desfazer essa confusão.
  \item \textbf{ENSO como ciclo regular}: é oscilador irregular
        (ruído $+$ não linearidade) — 3--7 anos, não relógio.
\end{itemize}

\subsection*{Demonstração de baixo custo}

A bifurcação do Caso~1 desenhada no quadro com setas de estabilidade
(análise gráfica de EDO — os alunos fazem sozinhos). Vídeo do tanque
giratório de convecção ou da série de TSM do Niño-3.4 (ondas reais de
El Niño/La Niña na tela).

\subsection*{Problemas sugeridos do livro (exercícios do Cap.~21)}

Aquecimento: resposta de Planck com números; períodos orbitais.
Aprofundamento: ganho com múltiplos feedbacks; equilíbrios do EBM com
albedo degrau. Síntese: ``por que o clima é previsível se o tempo não
é?''. \emph{Confira a numeração na sua cópia (v1.02b).}

\subsection*{Uso do notebook \texttt{cap21\_clima.ipynb}}

\begin{center}
\begin{tabular}{lll}
\hline
Caso & Conteúdo & Momento sugerido \\
\hline
1 & bifurcação gelo-albedo (bola de neve) & Aula 1 \\
2 & ciclos de Milankovitch e $Q_{65N}$ & Aula 1 \\
3 & Daisyworld completo & Aula 2 \\
4 & oscilador de recarga do ENSO & Aula 2 \\
\hline
\end{tabular}
\end{center}

Os quatro casos são papers clássicos (Budyko/Sellers 1969; Milankovitch;
Watson--Lovelock 1983; Jin 1997) em miniatura — vale dizer isso à
turma: eles estão rodando a literatura.

\subsection*{Exercícios complementares e soluções}

Nas notas: T1--T5 e N1--N4. Soluções em \texttt{cap21\_solucoes.ipynb}
(\emph{não distribuir}): T2 e T5 com \texttt{sympy}; N1 e N2 são os
mais ricos (histerese; retificação de Milankovitch). Pares:
T1$\leftrightarrow$N1, T3$\leftrightarrow$N2, T4$\leftrightarrow$N3,
T5$\leftrightarrow$N4.

\subsection*{Conexões com o resto do curso}

Balanço radiativo (Cap.~2); vapor e Clausius--Clapeyron (Cap.~4);
nuvens como feedback incerto (Caps.~6--7); circulação de Walker e
alísios (Cap.~11); caos e previsibilidade estatística (Cap.~20); o
Cap.~23 (módulo extra) aplica TUDO isto à perturbação antrópica; o Cap.~22 fecha o livro com a óptica atmosférica.
